\subsection{Curvature derivatives and exact Jacobi operators}
\label{app:rigid-differential}
Holding $q$ fixed while changing $u$ moves the physical point $x=e^{-u/2}q$. The curvature derivative must therefore include the induced point displacement. We first compute this total derivative, including its higher-order terms, and then analyze how tangent and base-point perturbations pass through the exponential map.
\begin{definition}[Euler operators]
\begin{align}
 \mathcal E_g&=x_g^\top\nabla_{x_g},&
 \mathcal D_g&=\partial_{u_g}-\tfrac12\mathcal E_g,\\
 X(q,u)_g&=e^{-u_g/2}q_g,&
 [\mathcal D_g,\mathcal D_h]&=0.
\end{align}
\end{definition}
\begin{proposition}[All-order curvature differentiation]
For $\widetilde{\mathcal J}\in C^r$, $|\alpha|\le r$,
\begin{align}\label{eq:total-curv-gradient}
 \partial_{u_g}F
 &=\partial_{u_g}\widetilde{\mathcal J}
   -\tfrac12\nabla_{x_g}\widetilde{\mathcal J}^{\top}x_g,\\
 \partial_u^\alpha F(q,u)
 &=\left(\prod_g\mathcal D_g^{\alpha_g}\widetilde{\mathcal J}\right)(X(q,u),u),\\
 D_u^jX[h^{(1)},\ldots,h^{(j)}]_g
 &=(-\tfrac12)^j\left(\prod_{i=1}^jh_g^{(i)}\right)x_g.
\end{align}
In particular,
\begin{align}
 F_{u_gu_g}
 &=\widetilde{\mathcal J}_{u_gu_g}
 -x_g^\top\widetilde{\mathcal J}_{x_gu_g}
 +\tfrac14x_g^\top\widetilde{\mathcal J}_{x_gx_g}x_g
 +\tfrac14\widetilde{\mathcal J}_{x_g}^\top x_g,\\
 F_{u_gu_h}
 &=\widetilde{\mathcal J}_{u_gu_h}
 -\tfrac12x_g^\top\widetilde{\mathcal J}_{x_gu_h}
 -\tfrac12x_h^\top\widetilde{\mathcal J}_{x_hu_g}
 +\tfrac14x_g^\top\widetilde{\mathcal J}_{x_gx_h}x_h
 \quad(g\ne h).
\end{align}
\end{proposition}
\begin{proof}
For $h,k\in\mathbb R^G$, put $V_h=D_uX[h]$ and $W_{h,k}=D_u^2X[h,k]$.
\begin{align}
 D_uF[h]&=\widetilde{\mathcal J}_u[h]+\widetilde{\mathcal J}_x[V_h],\\
 D_u^2F[h,k]
 &=\widetilde{\mathcal J}_{uu}[h,k]
 +\widetilde{\mathcal J}_{ux}[h,V_k]
 +\widetilde{\mathcal J}_{xu}[V_h,k]\\
 &\quad+\widetilde{\mathcal J}_{xx}[V_h,V_k]
 +\widetilde{\mathcal J}_x[W_{h,k}],\\
 \partial_{u_h}\bigl(\mathcal D_g\widetilde{\mathcal J}\bigr)(X(q,u),u)
 &=(\mathcal D_h\mathcal D_g\widetilde{\mathcal J})(X(q,u),u).
\end{align}
Induction on $|\alpha|$ gives the stated identity.
\end{proof}

\begin{proposition}[Softplus coordinate and second-order remainder]
\label{prop:softplus-rigid}
\begin{align}
 K(c)&=K_0+\log(1+e^c),&\psi(c)&=\log K(c),&s(c)&=(1+e^{-c})^{-1},\\
 \psi'(c)&=s(c)/K(c),&
 \psi''(c)&=\frac{s(c)(1-s(c))}{K(c)}-\frac{s(c)^2}{K(c)^2},\\
 \partial_cF&=\psi'(c)\partial_uF,&
 \partial_{cc}F&=\psi'(c)^2F_{uu}+\psi''(c)F_u.
\end{align}
If $c\in[c_-,c_+]$, $K(c)\in[K_-,K_+]$, then
\begin{align}
 b_-&:=\frac{s(c_-)^2}{K_+^2}
 \le\psi'(c)^2\le\frac{s(c_+)^2}{K_-^2}=:b_+,\\
 |\psi''(c)|&\le B_\psi:=\frac1{4K_-}+\frac1{K_-^2}.
\end{align}
For $c^+=c-\rho\psi'(c)\widehat b$,
\begin{align}
 u^+&=u-\rho\psi'(c)^2\widehat b+R_c,\\
 |R_c|&\le\tfrac12B_\psi\rho^2\psi'(c)^2|\widehat b|^2.
\end{align}
\end{proposition}
\begin{proof}
\begin{align}
 K'(c)&=s(c),&K''(c)&=s(c)(1-s(c)),\\
 \psi(c+v)&=\psi(c)+\psi'(c)v+
 v^2\int_0^1(1-r)\psi''(c+rv)\,dr.
\end{align}
The bounds apply on the full segment $[c,c+v]$.
\end{proof}

The softplus coordinate changes the effective curvature learning rate through $\psi'(c)^2$. The remainder is second order in the proposed coordinate step, but its cumulative contribution still has to be included in a finite-time budget. The next calculation isolates a different source of sensitivity: the transverse expansion of neighboring hyperbolic geodesics.

The Jacobi equations below are the constant-curvature model equations underlying Rauch comparison; see \citet{lee2018riemannian}. Using their exact solutions, rather than replacing them immediately by a curvature-radius upper bound, will allow the normalization and exponential differentials to cancel quantitatively.

\begin{lemma}[Jacobi initial-value equations]
\label{lem:rigid-jacobi-ivp}
For $\gamma(s)=\Exp_x^K(sv)$, $r=\|v\|$, $\lambda=\sqrt Kr$,
\begin{align}
 D_s^2J+R(J,\dot\gamma)\dot\gamma&=0,\\
 D_s^2J_\parallel&=0,&D_s^2J_\perp-\lambda^2J_\perp&=0.
\end{align}
For parallel identification $P_s=P_{x\to\gamma(s)}^K$,
\begin{align}
 J(0)=0,\ D_sJ(0)=w
 &\Longrightarrow
 J(s)=P_s\left[sw_\parallel+\frac{\sinh(\lambda s)}\lambda w_\perp\right],\\
 J(0)=a,\ D_sJ(0)=0
 &\Longrightarrow
 J(s)=P_s\left[a_\parallel+\cosh(\lambda s)a_\perp\right].
\end{align}
\end{lemma}
\begin{proof}
\begin{align}
 R(J,\dot\gamma)\dot\gamma
 &=-K(r^2J-\langle J,\dot\gamma\rangle\dot\gamma),\\
 (P_s^{-1}J)_\parallel''&=0,&
 (P_s^{-1}J)_\perp''&=\lambda^2(P_s^{-1}J)_\perp,\\
 j_\parallel(s)&=a_\parallel+sb_\parallel,&
 j_\perp(s)&=\cosh(\lambda s)a_\perp+
 \lambda^{-1}\sinh(\lambda s)b_\perp.
\end{align}
\end{proof}

\begin{lemma}[Exact exponential differential]
\label{lem:exact-jacobi}
For $s(z)=\sinh(z)/z$, $s(0)=1$,
\begin{align}
 D_v\Exp_x^K(v)&=P_1(\mathsf P_\parallel+s(\lambda)\mathsf P_\perp),\\
 \sigma(D_v\Exp_x^K(v))&=\{1,s(\lambda),\ldots,s(\lambda)\},\\
 \|D_v\Exp_x^K(v)\|_{\rm op}&=s(\lambda)\quad(d\ge2),\\
 \kappa(D_v\Exp_x^K(v))&=s(\lambda)\quad(d\ge2),\\
 \kappa(D_v\Exp_x^K(v))&=1\quad(d=1).
\end{align}
For a horizontal base-point variation,
\begin{align}
 D_x^{\rm hor}\Exp_x^K(v)
 &=P_1(\mathsf P_\parallel+\cosh(\lambda)\mathsf P_\perp),\\
 \|D_x^{\rm hor}\Exp_x^K(v)\|_{\rm op}&\le\cosh(\lambda).
\end{align}
\end{lemma}
\begin{proof}
\begin{align}
 \|D_v\Exp_x^K(v)[w]\|^2
 &=\|w_\parallel\|^2+s(\lambda)^2\|w_\perp\|^2,\\
 \|D_x^{\rm hor}\Exp_x^K(v)[a]\|^2
 &=\|a_\parallel\|^2+\cosh(\lambda)^2\|a_\perp\|^2.
\end{align}
\end{proof}

\begin{corollary}[Finite velocity perturbations]
For $v,w\in T_xH_K$, $\max(\|v\|,\|w\|)\le R$,
\begin{align}\label{eq:velocity-lip}
 d_K(\Exp_x^Kv,\Exp_x^Kw)
 &\le s(\sqrt KR)\|v-w\|.
\end{align}
For $x,y$ joined by a geodesic and $\|v\|\le R$,
\begin{align}
 d_K(\Exp_x^Kv,\Exp_y^K(P_{x\to y}^Kv))
 &\le\cosh(\sqrt KR)d_K(x,y).
\end{align}
\end{corollary}
\begin{proof}
\begin{align}
 v_\theta&=(1-\theta)v+\theta w,&\|v_\theta\|&\le R,\\
 d_K(\Exp_xv,\Exp_xw)
 &\le\int_0^1\|D\Exp_x(v_\theta)[w-v]\|\,d\theta\\
 &\le s(\sqrt KR)\|v-w\|.
\end{align}
For the second inequality, parallel transport $v$ along the base geodesic and integrate the horizontal differential.
\end{proof}

\begin{proposition}[Normalization differential]
\label{prop:normalization-rigid}
In an orthonormal tangent frame, let $N_\gamma(v)=v/(1+\gamma\|v\|)$, $\gamma>0$. For $v\ne0$, $e=v/\|v\|$,
\begin{align}
 DN_\gamma(v)
 &=\frac1{1+\gamma\|v\|}(I-ee^\top)
   +\frac1{(1+\gamma\|v\|)^2}ee^\top,\\
 DN_\gamma(0)&=I,&\|DN_\gamma(v)\|_{\rm op}&\le1,\\
 \|N_\gamma(v)-N_\gamma(w)\|&\le\|v-w\|,\\
 \|N_\gamma(v)\|&\le\min(\|v\|,\gamma^{-1}),\\
 \sqrt K\left\|\frac{-\eta m}{1+\gamma\sqrt K\|m\|}\right\|
 &\le\eta/\gamma.\label{eq:step-cap}
\end{align}
If $a_-I\preceq A\preceq a_+I$ and $\|g\|\le G$,
\begin{align}\label{eq:rigid-alignment}
 \langle g,N_\gamma(Ag)\rangle
 &\ge a\|g\|^2,&a&:=\frac{a_-}{1+\gamma a_+G},\\
 \|N_\gamma(Ag)\|&\le a_+\|g\|.
\end{align}
For blockwise normalization the same inequalities hold when $A=\operatorname{diag}(A_j)$ respects those blocks.
If each normalization block has a scalar preconditioner $A_j=a_jI$, the stronger uniform constant is
\begin{align}\label{eq:sharp-scalar-alignment}
 \langle g,N_\gamma(Ag)\rangle
 &\ge a_{\rm sc}\|g\|^2,&
 a_{\rm sc}&:=\frac{a_-}{1+\gamma a_-G}
 \ge\frac{a_-}{1+\gamma a_+G}.
\end{align}
This constant is attained when $g$ has norm $G$ and is supported on a block with $a_j=a_-$.
\end{proposition}
\begin{proof}
\begin{align}
 D\left(\frac v{1+\gamma\|v\|}\right)[h]
 &=\frac h{1+\gamma\|v\|}
 -\frac{\gamma\langle v,h\rangle v}{\|v\|(1+\gamma\|v\|)^2},\\
 \langle g,N_\gamma(Ag)\rangle
 &=\frac{\langle g,Ag\rangle}{1+\gamma\|Ag\|}
 \ge\frac{a_-\|g\|^2}{1+\gamma a_+G}.
\end{align}
For the scalar blocks,
\begin{align}
 \langle g,N_\gamma(Ag)\rangle
 &=\sum_j\frac{a_j\|g_j\|^2}{1+\gamma a_j\|g_j\|}
 \ge\sum_j\frac{a_j\|g_j\|^2}{1+\gamma a_jG}
 \ge\frac{a_-}{1+\gamma a_-G}\sum_j\|g_j\|^2,
\end{align}
since $a\mapsto a/(1+\gamma aG)$ is increasing.
\end{proof}

\begin{proposition}[Baillon--Haddad structure of normalization]
\label{prop:sharp-normalization-cocoercive}
For $\gamma>0$, define the convex radial potential
\begin{align}
 V_\gamma(v)&=\frac{\gamma\|v\|-\log(1+\gamma\|v\|)}{\gamma^2},&
 \nabla V_\gamma(v)&=N_\gamma(v).
\end{align}
Then normalization is firmly nonexpansive:
\begin{align}\label{eq:sharp-cocoercivity}
 \|N_\gamma(v)-N_\gamma(w)\|^2
 &\le\langle N_\gamma(v)-N_\gamma(w),v-w\rangle,\\
 \|v\|,\|w\|\le R
 &\Longrightarrow
 \langle N_\gamma(v)-N_\gamma(w),v-w\rangle
 \ge\frac{\|v-w\|^2}{(1+\gamma R)^2}.
\end{align}
The same statements hold for a product of separately normalized blocks.
\end{proposition}
\begin{proof}
The preceding differential gives $0\preceq\nabla^2V_\gamma\preceq I$. The Baillon--Haddad theorem \citep[Theorem~1.1]{bauschke2010baillon} therefore gives the first inequality. On the radius-$R$ ball, $\nabla^2V_\gamma\succeq(1+\gamma R)^{-2}I$; integration along the line segment gives the second. Sum the potentials for the blockwise statement.
\end{proof}

\begin{theorem}[Sharp sensitivity of a normalized geodesic step]
\label{thm:sharp-normalized-exp}
Fix $x\in H_K^d$, $d\ge2$, $\eta,\gamma>0$, and set
\begin{align}
 \mathcal S_x(m)&:=\Exp_x^K\!\left(\frac{-\eta m}{1+\gamma\sqrt K\|m\|}\right),&
 \alpha&:=\eta/\gamma,\\
 z&:=\frac{\gamma\sqrt K\|m\|}{1+\gamma\sqrt K\|m\|},&
 \Lambda(\alpha)&:=\max_{0\le z\le1}(1-z)\frac{\sinh(\alpha z)}{\alpha z}.
\end{align}
Use the continuous extension at $z=0$. The radial and transverse singular values of $D_m\mathcal S_x$ are exactly
\begin{align}\label{eq:sharp-step-singular-values}
 \sigma_\parallel&=\eta(1-z)^2,&
 \sigma_\perp&=\eta(1-z)\frac{\sinh(\alpha z)}{\alpha z}.
\end{align}
Consequently, the optimal global Lipschitz constant from the tangent metric to geodesic distance is $\eta\Lambda(\alpha)$, independently of $K$. In particular,
\begin{align}
 d_K(\mathcal S_x(m),\mathcal S_x(m'))
 &\le\eta\|m-m'\|\qquad\text{if }\alpha\le\alpha_c,\label{eq:sharp-step-lipschitz}\\
 \alpha_c&:=\inf_{\lambda>0}
 \frac{\lambda\sinh\lambda}{\sinh\lambda-\lambda}
 \ge\sqrt6.
\end{align}
For $d=1$, the optimal constant is $\eta$ for every $\alpha$.
\end{theorem}
\begin{proof}
The radial and transverse eigenspaces of the normalization differential agree with those of the Jacobi operator. Their singular values therefore multiply exactly, giving Eq.~\eqref{eq:sharp-step-singular-values}. Integrating the differential along the segment from $m$ to $m'$ proves the global upper bound. Conversely, small increments in a maximizing transverse direction approach the differential norm, so the constant is optimal.

Since the radial supremum is $\eta$ and the transverse factor equals one at $z=0$, $\Lambda(\alpha)=1$ precisely when
\begin{align}
 (1-\lambda/\alpha)\frac{\sinh\lambda}{\lambda}&\le1
 \quad(0<\lambda<\alpha)\\
 &\Longleftrightarrow\quad
 \alpha\le\frac{\lambda\sinh\lambda}{\sinh\lambda-\lambda}
 \quad(0<\lambda<\alpha).
\end{align}
The right-hand ratio is larger than $\lambda$. Thus taking its infimum over all positive $\lambda$ gives exactly the threshold $\alpha_c$.
Euler's product for $\sinh$ \citep[Eq.~4.36.1]{nistDLMF} gives the explicit sufficient threshold:
\begin{align}
 \log\frac{\sinh t}{t}
 &=\sum_{k\ge1}\log\left(1+\frac{t^2}{\pi^2k^2}\right)
 \le\frac{t^2}{6},\\
 \log\left[(1-z)\frac{\sinh(\alpha z)}{\alpha z}\right]
 &\le\log(1-z)+\alpha^2z^2/6
 \le-z+z^2\le0\qquad(\alpha^2\le6).
\end{align}
In dimension one only the radial singular value is present.
\end{proof}

The cancellation matters: multiplying a worst-case Jacobi expansion by a separate nonexpansiveness bound would give $\eta\sinh(\alpha)/\alpha$. Retaining their common radial variable gives the exact smaller factor $\eta\Lambda(\alpha)$, and in the stated regime removes this extra velocity-amplification factor entirely. Base-point sensitivity remains a separate issue in the coupled stability analysis.

The normalization map is nonexpansive and caps step length, but its radial derivative is smaller than its transverse derivative. Thus its benefit is a bound on finite-displacement amplification, with a corresponding change in the descent scale. Because normalization is nonlinear, an unbiased tangent oracle does not in general remain unbiased after normalization; the following example and bias bound make this distinction explicit.

\begin{proposition}[Normalization bias]
\begin{align}
 \Pr(\widehat g=2)&=1/3,&\Pr(\widehat g=-1)&=2/3,&\mathbb E\widehat g&=0,\\
 \mathbb EN_\gamma(\widehat g)
 &=\frac2{3(1+2\gamma)}-\frac2{3(1+\gamma)}
 =\frac{-2\gamma}{3(1+2\gamma)(1+\gamma)}\ne0.
\end{align}
\end{proposition}

\begin{proposition}[Optimal Lipschitz constant of the normalization Jacobian]
\label{prop:sharp-normalization-jacobian}
For $\gamma\ge0$,
\begin{align}
 \|DN_\gamma(v)-DN_\gamma(w)\|_{\rm op}&\le2\gamma\|v-w\|,\\
 \|N_\gamma(v+\xi)-N_\gamma(v)-DN_\gamma(v)\xi\|
 &\le\gamma\|\xi\|^2.
\end{align}
For $\gamma>0$, neither universal constant can be reduced.
\end{proposition}
\begin{proof}
For $v\ne0$, put $r=\|v\|$, $e=v/r$, and $a(r)=(1+\gamma r)^{-1}$. For a unit vector $h$, let $t=\langle e,h\rangle^2$ and $c=(1-\gamma r)/(1+\gamma r)\in[-1,1]$.
\begin{align}
 \|D^2N_\gamma(v)[h,h]\|^2
 &=\frac{\gamma^2}{(1+\gamma r)^4}
 \bigl[(1+ct)^2+4t(1-t)\bigr]\\
 &\le\frac{\gamma^2}{(1+\gamma r)^4}(1+6t-3t^2)
 \le4\gamma^2.
\end{align}
Polarization of the symmetric bilinear map $B=D^2N_\gamma(v)$ yields, for unit $h,k$,
\begin{align}
 \|B[h,k]\|
 &\le\tfrac14\bigl(\|B[h+k,h+k]\|+\|B[h-k,h-k]\|\bigr)\\
 &\le\tfrac{2\gamma}{4}(\|h+k\|^2+\|h-k\|^2)=2\gamma.
\end{align}
Integrate along the segment from $v$ to $w$, splitting at zero if necessary; $DN_\gamma$ is continuous there. The Taylor remainder is then bounded by
\begin{align}
 \left\|\int_0^1[DN_\gamma(v+s\xi)-DN_\gamma(v)]\xi\,ds\right\|
 &\le\int_0^12\gamma s\|\xi\|^2\,ds=\gamma\|\xi\|^2.
\end{align}
For a unit $e$, the limits
\begin{align}
 \lim_{r\downarrow0}\frac{\|DN_\gamma(re)-DN_\gamma(0)\|_{\rm op}}r&=2\gamma,\\
 \lim_{r\downarrow0}\frac{\|N_\gamma(re)-re\|}{r^2}&=\gamma
\end{align}
prove optimality. The case $\gamma=0$ is immediate.
\end{proof}

\begin{corollary}[Conditional normalization bias]
For predictable $v$ and $\mathbb E_t\xi=0$,
\begin{align}
 \|\mathbb E_tN_\gamma(v+\xi)-N_\gamma(v)\|
 &\le\gamma\mathbb E_t\|\xi\|^2,\\
 \mathbb E_t\|N_\gamma(v+\xi)-N_\gamma(v)\|^2
 &\le\mathbb E_t\|\xi\|^2
 -\mathbb E_t\|\xi-N_\gamma(v+\xi)+N_\gamma(v)\|^2
 \le\mathbb E_t\|\xi\|^2.
\end{align}
In the current momentum update,
\begin{align}
 \bar p_t&=\beta P_{t-1}p_{t-1}+(1-\beta)A_tg_t,\\
 \xi_t^p&=(1-\beta)A_t(\widehat g_t-g_t),\\
 \|\mathbb E_tN_\gamma(p_t)-N_\gamma(\bar p_t)\|
 &\le\frac{\gamma(1-\beta)^2a_+^2\sigma_q^2}{\mathfrak b_t}.
\end{align}
\end{corollary}
\begin{proof}
The linear Taylor term has conditional expectation $DN_\gamma(v)\mathbb E_t\xi=0$. Apply the optimal remainder bound. For the second moment, firm nonexpansiveness gives
\begin{align}
 \|N_\gamma(v+\xi)-N_\gamma(v)\|^2
 +\|\xi-N_\gamma(v+\xi)+N_\gamma(v)\|^2\le\|\xi\|^2.
\end{align}
Taking conditional expectations retains the nonnegative damping term rather than discarding it.
The bias constant is also optimal under conditional centering: in one dimension take $v=r>0$ and $\xi=\pm r^2$ with equal probabilities. Since $N_\gamma''(s)=-2\gamma/(1+\gamma s)^3$ for $s>0$, the centered second difference gives
\begin{align}
 \lim_{r\downarrow0}
 \frac{|\mathbb E N_\gamma(r+\xi)-N_\gamma(r)|}{\mathbb E\xi^2}=\gamma.
\end{align}
\end{proof}
